\documentclass[11pt,a4paper]{article}
\usepackage[margin=24mm]{geometry}
\usepackage{amsmath,amssymb,amsthm,booktabs,longtable,microtype}
\setlength{\emergencystretch}{2em}
\newtheorem{theorem}{Theorem}
\newcommand{\GL}{\operatorname{GL}}
\newcommand{\SL}{\operatorname{SL}}
\newcommand{\PGL}{\operatorname{PGL}}
\newcommand{\PSL}{\operatorname{PSL}}
\title{Linear groups and regular polytopes:\\the complete thesis classification}
\author{Human Compression project: the N012 principal recovery}
\date{}
\begin{document}\maketitle
This paper assembles the principal results of I.~A.~Markovskaya's
2025 thesis [T], in its fixed eighty-page edition. We acquire the
complete companion proofs [I,U,L,Z,H,Q], rather than repeat their
matrix calculations. One structural clause in Corollary~4.2.1 of
[T] requires correction; its polytope conclusion remains valid.
The classification below includes all twelve numbered principal
theorems and corollaries and both six-row tables of [T].

Every dimension is at least two. An admissible group is generated
by three involutions with a commuting pair, with the source convention
allowing repetitions. A rank-three regular-polytope group instead
requires an actual string C-marking, including the intersection
property and distinct nonidentity generators. For odd $q$ write
\[
 H_n(q)=\{g\in\GL_n(q):\det g\in\{1,-1\}\},\qquad
 Q_n(q)=H_n(q)/Z(\SL_n(q)).
\]
Thus $Q_n(q)$ is the source's $\PGL_n^{\pm1}(q)$.
It need not equal the quotient by the full centre of $H_n(q)$.

\section{The principal classifications}
\begin{theorem}
For $n\ge2$, the integral groups satisfy
\[
 \GL_n(\mathbb Z)\text{ admissible}\iff n\ge5,\qquad
 \PGL_n(\mathbb Z)\text{ admissible}\iff n=2\text{ or }n\ge5.
\]
Over a finite field the following table gives exactly the admissible
parameters. The regular-polytope parameters are the same except
that $\GL_2(2)$ and $\PGL_2(2)$ are omitted.
\begin{center}\small
\begin{tabular}{ll}\toprule
group&admissible parameters\\\midrule
$\GL_n(q)$&$q=2,\ n=2\text{ or }n\ge5$; or $q=3,\ n\ge5$\\
$\PGL_n(q)$&$n=2$; or $n\ge4,\ \gcd(n,q-1)=2$;\\
&or $n\ge5,\ \gcd(n,q-1)=1$\\
$H_n(q)$, $q$ odd&$n\ge5$\\
$Q_n(q)$, $q$ odd&$n=2,\ q\ne9$; or $n\ge4$\\\bottomrule
\end{tabular}
\end{center}
For odd $q=p^m$ and $n\ge5$, $H_n(q)$ has a rank-three regular
polytope of type $[4p,n]$ or $[4p,2n]$.
\end{theorem}
\begin{proof}
The full integral argument is [I]. It gives unit-root words for
every $n\ge5$, treats both degree-two groups directly, and excludes
degrees three and four by the explicitly retained finite quotients.
It is precisely [T, Theorem~2.2.1].

The two finite general-linear admissibility classifications are the
theorems labelled \texttt{gl} and \texttt{pgl} in [L]. Their parameter
sets agree exactly with [T, Theorems~3.2.1 and 3.3.1], and the two
tables in [L] recover every row of [T, Tables~1 and 2]. The integral
premise formerly retained by [L] is now supplied by [I]. This is
acyclic: [I] uses the separate exclusions $\PSL_3(2),\PSL_4(2)$,
not the finite general-linear classification that subsequently
uses integral reduction.

For $q=3$, $\GL_n(3)=H_n(3)$, so the actual marking and
intersection proof of [H] prove [T, Theorem~3.2.2] in every positive
dimension. For $q=2$, the positive groups are simple and the
explicitly retained admissibility and intersection criteria apply.
The remaining admissible group $S_3=\GL_2(2)$ has no distinct
commuting involutions, so it has no rank-three C-marking. The
projective polytope statement is proved in [L, Section~5], with
$\PGL_2(3)=S_4$ handled separately from the almost-simple cases;
the same $S_3$ exclusion completes [T, Corollary~3.3.1]. These
deductions are written out in [Q].

The first theorem of [H] gives the complete $H$ classification and
types, namely [T, Theorems~4.2.1 and 4.2.2]. It acquires the full
Scott obstruction of [Z] in degrees two, three and four. The second
theorem of [H] supplies actual C-markings for all scalar quotients
in dimensions at least five. In dimension four [Z] proves the
$Q$ assertion over every odd field, with both centre cases treated.
Finally [Q] proves the precise small quotient cases: $Q_2(q)$ is
positive exactly when $q\ne9$, and every $Q_3(q)$ is negative.
Its central direct products use the actual perfect-group lift of
[Z]. This proves [T, Theorem~4.3.1 and Corollary~4.3.1], including
their negative assertions.
\end{proof}

\section{The three quotient corollaries}
Put $Z_1=Z(H_n(q))$, $Z_2=Z(\SL_n(q))$,
$d=\gcd(n,q-1)$ and $e=(q-1)/d$.
\begin{theorem}
For odd $q=p^m$ and $n\ge5$, the following conclusions hold.
\begin{enumerate}
\item If $n$ is even, the full-centre quotient $H_n(q)/Z_1$
is $\PSL_n(q)$ when $e$ is even. When $e$ is odd, it is the
index-two extension $\PSL_n(q)\rtimes C_2$ embedded in
$\PGL_n(q)$, with complement represented by
$\operatorname{diag}(-1,1,\ldots,1)$. In both cases it has a
rank-three regular polytope of type $[4p,n]$.
\item If $n$ is odd, $H_n(q)/Z_1\cong\PSL_n(q)$ has a
rank-three regular polytope of type $[4p,n]$.
\item If $n$ is odd, $H_n(q)/Z_2\cong\PSL_n(q)\times C_2$
has a rank-three regular polytope of type $[4p,n]$ or $[4p,2n]$.
\end{enumerate}
The first clause corrects the structural part of [T,
Corollary~4.2.1]; the second and third recover Corollaries~4.2.2
and 4.2.3 as stated.
\end{theorem}
\begin{proof}
The centre theorem of [Z] proves exactly these group isomorphisms.
When $e$ is even, $|Z_1:Z_2|=2$, and otherwise $Z_1=Z_2$.
Odd $n$ always has even $e$, since $-I$ has determinant $-1$.
The second theorem of [H], applied first to $Z_1$ and then to
$Z_2$, proves the asserted types by an actual quotient intersection
calculation. In its notation $|AC|=4p$ and
$|\langle A,C\rangle|=8p$; this dihedral subgroup meets the
scalar kernel trivially. The monomial product $BC$ has order
$n$ or $2n$, and has order exactly $n$ after dividing by the
full centre. No cyclic-normal-subgroup criterion is applied to
the direct product in the third clause.
\end{proof}

For clarity, the printed even-dimension assertion $Z_1=Z_2$
on [T, p.~69] is false: at $n=6,q=5$, $2I$ has determinant
$-1$, $|Z_1|=4$ and $|Z_2|=2$. Thus the structural clause
of its numbered corollary genuinely changes. Its type and existence
assertions do not change. The corrected classification is the
principal disposition recorded here.

\section{Attribution and the exact boundary}
All source principal results are attributed to Markovskaya and the
collaborators named in her thesis. Its pages 6--8 credit the integral
theorem jointly with the supervisor, the difficult five- and
six-dimensional triples to the author, the finite general-linear
negative cases jointly with the supervisor, and the explicit
positive constructions and Chapter~4 results to the author.
The companion papers distinguish the source matrices and Scott
method from the complete common root arguments, repaired identities,
centre calculation, central lift and uniform intersection proofs
supplied in this project. No general priority claim is made.

The precise primary ledger has twelve rows. Supporting lemmas and
survey references are recorded separately. In particular the
rank-two Dickson--Levchuk assertions are not premises of the new
root argument or the small quotient proof. The external inputs
actually retained are the stated simple-linear admissibility
classification, Sjerve--Cherkassoff's $\PGL_2$ theorem,
Conder--Oliveros and the usual group--polytope correspondence.
Their full external proof and software costs remain unknown.

The corrections in [Z,H] are retained explicitly: the printed
degree-four matrix on [T, p.~72] is not a projective involution,
and the seven-dimensional word on p.~67 needs the recorded square.
The former defect was already found in the project's phase~22
audit. The replacement uses corrected witnesses. For the finite
$\GL_6(3)$ polytope it does not claim that its marking is the
original integral marking with the source's pair orders $24,12$.
It proves the same principal existence conclusion.

The assembly uses no new finite search, no new external
classification and no new Lean declaration. Complete companion
manuscripts and their exact validation and distinct review records
are required dependencies. This paper proposes one corrected
principal recovery of N012; independent whole-source review and
publication acceptance remain separate. It asserts no new
open-conjecture resolution and no compression factor.

\begin{thebibliography}{9}\small
\bibitem{T}[T] I.~A.~Markovskaya, \emph{Linear groups over finite
fields and over the ring of integers as automorphism groups of
regular polytopes of rank 3}, thesis, Siberian Federal University,
2025, fixed eighty-page edition; supervisor Ya.~N.~Nuzhin.
\bibitem{I}[I] \emph{Integral linear groups generated by three involutions}.
\bibitem{U}[U] \emph{Root parameters and fixed-space obstructions}.
\bibitem{L}[L] \emph{Three involutions in finite general linear groups}.
\bibitem{Z}[Z] \emph{Central quotients of the determinant-sign subgroup}.
\bibitem{H}[H] \emph{Involutions and intersections in groups with
determinant $\pm1$}, complete accepted second edition; the portable
copy clarifies the product and dihedral subgroup orders explicitly.
\bibitem{Q}[Q] \emph{The remaining quotients and the linear polytope
classifications}.
\end{thebibliography}
\end{document}
